%
% Abstract
%
\chapter{Abstract}
\begin{SingleSpace}
Urban transportation systems produce heterogeneous traffic and sensor data that must be cleaned, validated, and prepared before analysis. This thesis presents \textbf{Flowmatic}, an intelligent preparation and inference platform for batch uploads and smart-city pipelines, evaluated on a semi-synthetic Astana case study and public Hugging Face telemetry bundles.

Flowmatic implements six-dimensional data quality assessment, automated cleaning, and export in a NestJS/Vue Docker stack. A measured batch experiment ingests 30\,000 Astana records in 406\,ms with quality score 100/100. A smart-city workbench connects an Astana geospatial simulator, Manual/Auto model routing over a documented registry, and medallion lake export.

Eight neural checkpoints---PatchTST, iTransformer, TranAD, TimesBlock, STGCN, SAITS, a Transformer severity classifier, and DLinear---are trained under a shared temporal hold-out protocol (70/15/15 split, three seeds). The severity classifier reaches macro-F1 = $0.911 \pm 0.017$ on the test partition; TranAD reconstruction RMSE is $0.035 \pm 0.004$ on Astana windows. Weights and reproduction scripts are published on Hugging Face (\texttt{@pushthetempo/flowmatic-*}).

The dissertation demonstrates that preparation, multimodal inference, and operational deployment can be delivered as one auditable platform with reproducible artifacts.

\end{SingleSpace}
\keywords{Data Preparation, Intelligent Transportation Systems, Smart Cities, Model Routing, Time Series, Anomaly Detection, Flowmatic}
